\documentclass[11pt,a4paper]{article}
\usepackage[margin=25mm]{geometry}
\usepackage{amsmath,amssymb,amsthm,booktabs,array}
\usepackage[T1]{fontenc}
\usepackage{lmodern}
\usepackage[colorlinks=true,linkcolor=blue,citecolor=blue,urlcolor=blue]{hyperref}
\newtheorem{theorem}{Theorem}
\newtheorem{corollary}[theorem]{Corollary}
\DeclareMathOperator{\diag}{diag}
\title{Extending Towse--Campbell's rank-one constructions:\\
a counterexample to Finch's balanced order-six conjecture}
\author{Haobo Ma\textsuperscript{1} \and Wenlin Zhang\textsuperscript{2,1}}
\date{November 16, 2026}
\begin{document}
\maketitle
\begin{center}
\small
\textsuperscript{1}ChronoAI Pte Ltd, Singapore\\
\textsuperscript{2}National University of Singapore, Singapore
\end{center}
\begin{abstract}
Extending Towse and Campbell's integer rank-one constructions
(2016, Section 2, Theorem 1 and Corollaries 1--2) to the zerofree setting,
we refute Finch's Conjecture 4 in \emph{Small Matrices with Small Inverses:
Unimodular Zerofree Cases} (arXiv:2605.03691v4, June 12, 2026,
Section 5, p.~7) by constructing
a positive integral matrix of order six with determinant $-1$, no zero
entry in either it or its inverse, and maximum absolute entry $3$ in both.
A diagonal sign matrix
plus a rank-one matrix gives the construction. Every even order at least two admits such
a positive matrix with any prescribed common maximum at least three;
every odd order at least three admits one with common maximum three.
We also give positive order-six examples with maxima $(3,4)$ and $(3,5)$,
and identify the parts of Finch's extremal conjecture covered by his later bounds.
\end{abstract}

\section{A rank-one family}
For an integer matrix $M=(m_{ij})$ of order $n$, write
\[
 \alpha(M)=\max_{1\leq i,j\leq n}|m_{ij}|,
 \qquad
 \beta(M)=\max_{1\leq i,j\leq n}|(M^{-1})_{ij}|.
\]
Here $\beta(M)$ is used only when $M$ is invertible. Following Finch
\cite[pp.~1--2]{finch}, $M$ is \emph{unimodular} if
$\det M\in\{1,-1\}$ and \emph{zerofree} if both $M$ and $M^{-1}$
have no zero entries. A balanced matrix has $\alpha(M)=\beta(M)$.
Positivity below refers to the entries of the displayed matrix, without
selecting a canonical representative.
Finch's Conjecture 4 states: ``No $6\times6$ matrices exist that are
tied to $\alpha=\beta=3$.''
The following family supplies counterexamples without an enumeration
of equivalence classes.
Towse and Campbell \cite[Section 2, Theorem 1 and Corollaries 1--2]{towse}
record the integer rank-one determinant and inverse identities for
$P=I+\mathbf u\mathbf v^{\mathsf T}$, including the orthogonal
and involutory cases $\mathbf v^{\mathsf T}\mathbf u=0,-2$;
they attribute the orthogonal construction to Ortega.
The inverse identity is also the Sherman--Morrison formula; see
Higham \cite[``Sherman--Morrison Formula'']{higham}.
The diagonal-sign specialization below gives positivity, zerofreeness and equal
prescribed entry maxima, conclusions not stated in \cite{towse}.
The rank-one identities themselves are the cited prior method.
We give the determinant and inverse calculations in full.

Let $\mathbf e=(1,\ldots,1)^{\mathsf T}\in\mathbb Z^n$ and
$J=\mathbf e\mathbf e^{\mathsf T}$. For signs
$s_1,\ldots,s_n\in\{1,-1\}$ put
\[
 D=\diag(s_1,\ldots,s_n),\qquad
 \mathbf s=D\mathbf e,\qquad
 \delta=\sum_{i=1}^n s_i.
\]
Our family specializes Towse and Campbell's construction with
$\mathbf u=c\mathbf s$, $\mathbf v=\mathbf e$ and $A=DP$.
The inner product is $\mathbf v^{\mathsf T}\mathbf u=c\delta$,
so the two unimodular branches below correspond to $0$ and $-2$.

\begin{theorem}\label{thm:family}
Let $n\geq2$ and let $c\geq2$ be an integer. The positive integer
matrix $A=D+cJ$ is unimodular if and only if
\[
 \delta=0\quad\hbox{or}\quad(c,\delta)=(2,-1).
\]
In either case $A$ is zerofree and
$\alpha(A)=\beta(A)=c+1$. More explicitly,
\begin{align}
 \det A&=\det(D)(1+c\delta),\label{eq:det}\\
 A^{-1}&=D-\frac{c}{1+c\delta}\,
               \mathbf s\mathbf s^{\mathsf T}.\label{eq:inverse}
\end{align}
Formula \eqref{eq:inverse} also holds without unimodularity.
\end{theorem}

\begin{proof}
Since $D^2=I$, one has
$A=D(I+c\mathbf s\mathbf e^{\mathsf T})$.
Expand the determinant of the second factor by multilinearity in its
columns. A term using two columns from
$c\mathbf s\mathbf e^{\mathsf T}$ is zero because those columns
are proportional. The term using none is $1$, and the term using
only column $i$ is $cs_i$. This proves \eqref{eq:det}.
Because $c\geq2$ and $\delta$ is an integer, $1+c\delta\ne0$.
Moreover, $|\det A|=1$ is equivalent to
$1+c\delta\in\{1,-1\}$, that is, $c\delta=0$ or $c\delta=-2$.
These are precisely the two alternatives in the statement.

For the inverse, use
$D\mathbf s=\mathbf e$,
$\mathbf e^{\mathsf T}D=\mathbf s^{\mathsf T}$, and
$\mathbf e^{\mathsf T}\mathbf s=\delta$ to obtain
\[
 (D+c\mathbf e\mathbf e^{\mathsf T})
 \left(D-\frac{c}{1+c\delta}\mathbf s\mathbf s^{\mathsf T}\right)
 = I+\left(c-\frac{c+c^2\delta}{1+c\delta}\right)
               \mathbf e\mathbf s^{\mathsf T}=I.
\]
This proves \eqref{eq:inverse}.

The off-diagonal entries of $A$ equal $c$, and its diagonal entries
are $c+s_i\geq1$. Under either admissible alternative both signs
occur: $\delta=0$ requires equal positive and negative counts, while
$\delta=-1$ with $n\geq2$ requires odd $n\geq3$ and at least one
positive sign. Thus $\alpha(A)=c+1$.
If $\delta=0$, the off-diagonal entries of $A^{-1}$ have absolute
value $c$ and its diagonal entries are $s_i-c$, whose absolute
values are $c-1$ and $c+1$. If $(c,\delta)=(2,-1)$, then
$A^{-1}=D+2\mathbf s\mathbf s^{\mathsf T}$: its off-diagonal
entries have absolute value $2$ and its diagonal entries are
$s_i+2\in\{1,3\}$. Both inverses are integral and zerofree,
with the asserted maximum.
\end{proof}

\begin{corollary}\label{cor:orders}
For every even $n\geq2$ and every integer $k\geq3$, a positive
unimodular zerofree matrix exists with $\alpha=\beta=k$.
For every odd $n\geq3$, one exists with $\alpha=\beta=3$.
\end{corollary}
\begin{proof}
For $n=2m$, take $m$ positive and $m$ negative signs and $c=k-1$.
For $n=2m+1\geq3$, take $m$ positive and $m+1$ negative signs
and $c=2$. Apply Theorem~\ref{thm:family}.
\end{proof}

\section{The six-dimensional counterexample}
Take $D=\diag(1,1,1,-1,-1,-1)$ and $c=2$. Then
\begin{align*}
 A&=\begin{pmatrix}
 3&2&2&2&2&2\\
 2&3&2&2&2&2\\
 2&2&3&2&2&2\\
 2&2&2&1&2&2\\
 2&2&2&2&1&2\\
 2&2&2&2&2&1
 \end{pmatrix},&
 B&=\begin{pmatrix}
 -1&-2&-2&2&2&2\\
 -2&-1&-2&2&2&2\\
 -2&-2&-1&2&2&2\\
 2&2&2&-3&-2&-2\\
 2&2&2&-2&-3&-2\\
 2&2&2&-2&-2&-3
 \end{pmatrix}.
\end{align*}
Theorem~\ref{thm:family} gives $B=A^{-1}$ and
$\det A=(-1)^3=-1$. All entries of $A$ are in $\{1,2,3\}$;
all entries of $B$ are in $\{-3,-2,-1,2\}$. Both contain an
entry of absolute value $3$ and neither contains zero. Thus
Conjecture 4 is refuted; its nonexistence assertion is replaced by
the existence statement in Corollary~\ref{cor:orders}.

Finch identifies matrices up to $M\mapsto PMQ$, where $P,Q$ are
signed-permutation matrices \cite[p.~2]{finch}. Such multiplication
only permutes and changes signs of entries, and
$(PMQ)^{-1}=Q^{-1}M^{-1}P^{-1}$. Consequently unimodularity,
zerofreeness, $\alpha$ and $\beta$ are preserved. In particular,
the canonical representative of the orbit of $A$ is also a
counterexample; no choice of representative changes the conclusion.

\section{Two neighbouring six-dimensional cases}
Finch writes in \cite[Section 8, p.~11]{finch}:
``Results for the cases $(6,3,4)$ and $(6,3,5)$''.
Here the triples specify $(n,\alpha,\beta)$. The following examples supply
existence in both cases, rather than an enumeration of their orbits.

\begin{corollary}\label{cor:neighbours}
There are positive unimodular zerofree matrices of order six with
$(\alpha,\beta)=(3,4)$ and $(3,5)$, respectively.
\end{corollary}
\begin{proof}
Consider the following two pairs:
\begin{align*}
 C_4&=\begin{pmatrix}
 1&3&1&3&2&1\\1&2&1&1&1&1\\2&3&2&3&3&2\\
 1&1&3&2&2&2\\1&2&2&3&2&1\\1&2&2&3&3&2
 \end{pmatrix},&
 E_4&=\begin{pmatrix}
 1&-2&2&1&-1&-2\\-1&2&-1&-1&1&1\\-2&2&-1&-1&2&1\\
 3&-3&1&2&-2&-2\\-4&3&-1&-3&3&3\\4&-3&1&3&-4&-2
 \end{pmatrix},\\[4pt]
 C_5&=\begin{pmatrix}
 2&3&1&2&2&1\\1&2&1&3&2&1\\1&1&1&1&1&1\\
 2&2&1&1&3&1\\1&1&3&3&3&2\\1&1&3&2&2&2
 \end{pmatrix},&
 E_5&=\begin{pmatrix}
 3&-4&3&-2&5&-5\\-1&2&-2&1&-3&3\\2&-2&-1&-1&2&-1\\
 1&-1&1&-1&2&-2\\-1&1&-1&1&-1&1\\-4&4&1&2&-5&4
 \end{pmatrix}.
\end{align*}
Integer multiplication gives $C_jE_j=E_jC_j=I_6$ for $j=4,5$.
Taking determinants proves $\det C_j\in\{1,-1\}$, since both factors
are integers; in fact both determinants are $1$.
The displayed entries give $\alpha(C_j)=3$, $\alpha(E_j)=j$,
and show that neither matrix has a zero entry. Thus $E_j=C_j^{-1}$
has all the required properties.
\end{proof}

\section{Relation to Finch's other results}
Finch's related paper \cite[Sections 2--7]{roots} studies minimizing
$\max\{\alpha,\beta\}$ and lists explicit examples in small
orders. The construction above instead fixes the common maximum.
For example, the odd-order construction with $n=3$ gives Finch's
balanced matrix from \cite[Section 2, p.~3]{finch}, after moving
the positive sign to the third position. The formulas do not
enumerate signed-permutation orbits.
In particular, Corollary~\ref{cor:orders} supplies balanced height-three
examples in order ten and every larger order. This gives existence
beyond the examples in \cite[Addendum I, p.~15]{roots}, while leaving
the minimum-height and orbit-count questions separate.

Let $\gamma_n(k)$ be the largest inverse-entry maximum among unimodular
zerofree matrices of order $n$ with $\alpha=k$, as in
\cite[Section 9, p.~11]{finch}. The finite set of integer matrices with
entries of absolute value at most $k$ makes this a maximum whenever
the defining set is nonempty.
Corollary~\ref{cor:orders} immediately yields the following bounds.
\begin{corollary}\label{cor:gamma}
For $m\geq1$ and $k\geq3$ one has $\gamma_{2m}(k)\geq k$;
for $m\geq1$ one has $\gamma_{2m+1}(3)\geq3$.
\end{corollary}
\begin{proof}
The matrices in Corollary~\ref{cor:orders} belong to the defining sets
and have the indicated inverse-entry maxima.
\end{proof}

Two equalities conjectured in \cite[Conjecture 6, p.~11]{finch} are already
consequences of Finch's later bounds, first released July 8, 2026,
rather than new conclusions here:
\begin{equation}\label{eq:prior}
 \gamma_3(k)=k(2k-1),\qquad
 \gamma_4(k)=k(2k-1)(2k+1)\qquad(k\geq3).
\end{equation}
Indeed, \cite[Appendix 5.2, p.~17; Theorem (ii), p.~2]{large}
bounds the inverse-entry maximum by these expressions for all
unimodular matrices of the respective orders, even allowing zero entries.
The zerofree matrices $U,V$ in \cite[Section 9, p.~12]{finch}
attain the respective bounds for every $k\geq3$.
Restricting the upper bounds to the zerofree subclass and using these
lower-bound examples proves \eqref{eq:prior}.
The exceptional cases $\gamma_3(2)=5$ and $\gamma_4(2)=26$, and the
order-five and order-six parts of Conjecture 6, require separate arguments.

\begin{center}
\small
\begin{tabular}{@{}p{0.44\linewidth}p{0.51\linewidth}@{}}
\toprule
Source item & Coverage\\
\midrule
\cite[Section 2, Theorem 1; Corollaries 1--2]{towse}
& Prior proof cited: rank-one determinant and inverse identities,
including the zero and $-2$ inner-product branches. The zerofree
height conclusions are the specialization proved here.\\[3pt]
Conjecture 4, Section 5, p.~7 & Refuted and corrected by an explicit
matrix; Corollary~\ref{cor:orders} gives an even-order extension.\\[3pt]
Balanced examples in Sections 1--4, 6 & Not posed; proved: every order $n\geq2$
admits common maximum $3$; the order-three example is recovered.
Orbit counts are not covered.\\[3pt]
Conjectures 1, 2, 3, pp.~3--5 & Not covered: their hypotheses
require $\alpha=2$, unlike this family.\\[3pt]
Propositions 0 and 5, pp.~1, 9 & Prior proofs cited in \cite{finch};
the lower bound $\alpha,\beta\geq2$ and
the order-two count $2\varphi(k)-1$ are not
derived here.\\[3pt]
Section 8, p.~11, $(6,3,4)$, $(6,3,5)$ & Answered: yes for existence,
by Corollary~\ref{cor:neighbours}; orbit counts not covered.\\[3pt]
Section 7, p.~9; Section 9, p.~12 & Not covered: improvements to the
large order-seven or order-six cofactors requested there.\\[3pt]
Conjecture 6, Section 9, p.~11 & Prior proof cited: Finch \cite{large}
implies \eqref{eq:prior}, for $k\geq3$. Exceptional $k=2$ cases and
orders five and six not covered. Corollary~\ref{cor:gamma} gives
additional lower bounds, not extremal equalities.\\[3pt]
\cite[Sections 2--6; Addendum I, p.~15]{roots}; MSE \cite{mse}
& Prior constructions cited: small-height and order-nine examples.
Not posed; proved: balanced height-three existence in order ten and above.
Minimum height and complete orbit classification are not covered.\\
\bottomrule
\end{tabular}
\end{center}

\section{Verification scope}
The Lean file
\href{https://github.com/the-omega-institute/trureturing/blob/3ede1d0bfa6d3c063b82234699d279136b6f0c42/D5/S3/QuadraticForms/FinchBalancedZerofreeRefutation.lean}{\texttt{FinchBalancedZerofreeRefutation.lean}}
declares \texttt{result} in namespace\\
\texttt{Trureturing.OpenProblems.FinchUnimodularZerofree}.\\
It refutes \texttt{claim}, the integral order-six nonexistence assertion
with unit determinant, zerofreeness and entrywise maxima
$\alpha=\beta=3$, using $A,B$ and Mathlib's nonsingular inverse.
The all-order family and its corollaries have written proofs;
\eqref{eq:prior} follows from the cited results.
Corollary~\ref{cor:neighbours} gives finite integer certificates.
The accompanying \texttt{verify.py} checks all three inverse pairs and
finite family instances with exact rational arithmetic, with ranges
in its output. Lean coverage is the balanced order-six example alone.

\begin{thebibliography}{9}
\bibitem{finch}
Steven Finch, \emph{Small Matrices with Small Inverses:
Unimodular Zerofree Cases}, arXiv:2605.03691v4, June 12, 2026.
\url{https://arxiv.org/abs/2605.03691v4}.

\bibitem{roots}
Steven Finch, \emph{$n$th Roots of $n$th Powers},
arXiv:2602.17719v8, April 15, 2026.
\url{https://arxiv.org/abs/2602.17719v8}.

\bibitem{large}
Steven Finch, \emph{Small Matrices with Large Inverses:
Unimodular $4\times4$ Cases}, arXiv:2607.07688v4, August 17, 2026.
First version released July 8, 2026.
\url{https://arxiv.org/abs/2607.07688v4}.

\bibitem{mse}
S. Finch, \emph{Optimizing zerofree unimodular matrices},
Mathematics Stack Exchange, February 24, 2026;
answer by S. Finch, March 29, 2026.
\url{https://math.stackexchange.com/questions/5126078}.

\bibitem{higham}
Nick Higham, \emph{What Is the Sherman--Morrison--Woodbury Formula?},
September 29, 2020; updated November 12, 2020.
\url{https://nhigham.com/2020/09/29/what-is-the-sherman-morrison-woodbury-formula/}.

\bibitem{towse}
Christopher Towse and Eric Campbell,
\emph{Constructing Integer Matrices with Integer Eigenvalues},
author-hosted manuscript dated March 25, 2016,
Section 2, Theorem 1 and Corollaries 1--2, p.~2.
\url{https://ericthewry.github.io/pdfs/imies.pdf}.
\end{thebibliography}
\smallskip
\noindent\small
\copyright\ 2026 Haobo Ma and Wenlin Zhang.
Licensed under \href{https://creativecommons.org/licenses/by/4.0/}{CC BY 4.0}.
\end{document}
